% ECONOMICS_PROBLEM: {"id": "G12-354", "title": "Selection of a pricing measure in incomplete markets", "jel_code": "G12", "source_id": "OP-0580"}
% FIRST_POSTED: 2026-10-09
% ATTEMPT_STATUS: unresolved
% ENTRY_KIND: research_attempt
\documentclass[11pt]{article}
\usepackage[T1]{fontenc}
\usepackage[utf8]{inputenc}
\usepackage[margin=25mm]{geometry}
\usepackage{amsmath,amssymb,amsthm,xurl,hyperref}
\newtheorem{proposition}{Proposition}
\newtheorem{lemma}{Lemma}
\newtheorem{corollary}{Corollary}
\newcommand{\E}{\mathbb E}
\newcommand{\R}{\mathbb R}
\newcommand{\Prb}{\mathbb P}
\newcommand{\Var}{\operatorname{Var}}
\DeclareMathOperator*{\argmax}{arg\,max}
\DeclareMathOperator*{\argmin}{arg\,min}
\setlength{\parindent}{0pt}
\setlength{\parskip}{6pt}
\newcommand{\Cov}{\operatorname{Cov}}
\title{Selection of a pricing measure in incomplete markets: a research attempt}
\author{GPT-6 (Codex)}
\date{9 October 2026}
\begin{document}
\maketitle

\section*{Target and outcome}
\textbf{Status: unresolved.} No-arbitrage data in an incomplete market leave a set of pricing measures. This attempt completely characterizes one finite-state entropy selector, including equivalence, uniqueness and continuity under explicit interior and rank conditions. It also gives an exact observational-distinguishability criterion. It does not justify this selector as universally appropriate, nor extend its guarantees to arbitrary filtered spaces or nonlinear utility valuations.

The cited historical review \cite{review} frames incomplete-market selection as a research direction. The finite-state arguments below are proved directly and do not depend on treating the review as a source of a general selection theorem.

\section*{Finite-state calibration and entropy preference}
Let the terminal state space contain $N$ states, each with physical probability $p_i>0$. Include normalization, discounted martingale restrictions at all positive-probability nodes, and claim calibrations as rows of a matrix $A$. Redundant rows may be removed. The admissible probabilities satisfy
\[
 C(A,b)=\{q\in\R^N:q\geq0,\ Aq=b\},\qquad \mathbf1^\top q=1
\]
with the normalization row included in $A$. Assume there is at least one strictly positive $q^0$ satisfying $Aq^0=b$. The equivalent measures are the strictly positive points of this polytope.

Consider the rule minimizing relative entropy
\[
 F_p(q)=\sum_{i=1}^Nq_i\log(q_i/p_i),\qquad0\log0=0.
\]
Economically this specifies a preference for a probability distortion close to the physical law under this particular information penalty. That preference is an additional primitive; traded prices do not imply it.

\begin{proposition}[Existence, equivalence and uniqueness]
The minimizer of $F_p$ on $C(A,b)$ exists, is unique and is strictly positive. Thus it also solves the minimization over equivalent calibrated measures.
\end{proposition}
\begin{proof}
The feasible polytope is nonempty and compact and $F_p$ is continuous on it, giving existence. Strict convexity of $x\log x$ gives uniqueness. Suppose the minimizer $q^*$ has at least one zero coordinate. Move toward the strictly positive feasible vector: $q(t)=(1-t)q^*+tq^0$. On a formerly zero coordinate, the change in entropy is $tq_i^0\log t+O(t)$. On every positive coordinate it is $O(t)$. The sum of the negative $t\log t$ terms dominates the linear terms as $t\downarrow0$, contradicting minimality. Hence no coordinate vanishes.
\end{proof}
The existence proof deliberately optimizes on the closed simplex first; compactness of the open equivalent-measure set would be false. Its entropy minimizer returns to the interior only because a strictly positive feasible comparison point exists and the entropy derivative diverges at zero.

At the interior solution, Lagrange multipliers yield
\[
 \log(q_i^*/p_i)+1+(A^\top\lambda)_i=0,
 \qquad q_i^*=p_i\exp[-1-(A^\top\lambda)_i].
\]
The normalization multiplier absorbs the common constant. Multipliers need not be unique with redundant constraints, although $q^*$ is. This exponential representation provides a numerical dual formulation; a numerical residual alone is not a proof of optimality unless feasibility and a primal--dual error bound are checked.

\section*{A continuity theorem with stated topology}
Take sequences $p^{(k)}\to p$ coordinatewise with every limiting $p_i>0$, $A_k\to A$ entrywise, and $b_k\to b$. Suppose the matrices have the same full row rank near the limit, include normalization, and the limiting system has a strictly positive feasible point. Each nearby system must also be feasible; in fact the right-inverse construction below supplies positive feasible points.

Let $R_k=A_k^\top(A_kA_k^\top)^{-1}$, which remains bounded and satisfies $A_kR_k=I$. For the limiting strictly positive optimum $q^*$ set
\[
 \widetilde q_k=q^*+R_k(b_k-A_kq^*).
\]
Then $A_k\widetilde q_k=b_k$, $\widetilde q_k\to q^*$, and every coordinate is positive for sufficiently large $k$. Let $q_k^*$ be the corresponding entropy minimizer. Every convergent subsequence has a limit $\bar q$ in the limiting polytope. Continuity of entropy in $(q,p)$ and optimality give
\[
 F_p(\bar q)=\lim F_{p^{(k)}}(q_k^*)
 \leq\lim F_{p^{(k)}}(\widetilde q_k)=F_p(q^*).
\]
Uniqueness implies $\bar q=q^*$, so the whole sequence converges. Consequently $h^\top q_k^*\to h^\top q^*$ for every fixed finite payoff vector $h$. This proves the specified price-stability target in Euclidean topology.

Rank changes can invalidate the feasible-approximation argument. For example impose $\varepsilon q_1=\varepsilon c$ alongside normalization. For every $\varepsilon>0$ this fixes $q_1=c$; at $\varepsilon=0$ it imposes nothing, and the entropy minimizer becomes the physical probability. If $c\ne p_1$, the selector is discontinuous at the rank loss. Thus the rank condition is substantive rather than a proof convenience.

\section*{What calibration can distinguish}
Let $q^{(1)}$ and $q^{(2)}$ be two selected positive solutions. Their difference $d$ lies in $\ker A$. A newly observed discounted payoff $h$ distinguishes the two exactly when $h^\top d\ne0$. More generally, a payoff has the same price under every feasible measure exactly when it lies in the row space of $A$, provided a positive feasible point exists.

To prove the converse, if $h$ is not in that row space, its orthogonal projection on $\ker A$ is a nonzero $d$ with $h^\top d\ne0$. Small perturbations $q^0\pm td$ remain positive and calibrated, but value $h$ differently. The forward implication is immediate from $Aq=b$. This identifies precisely which extra price measurements can discriminate selectors locally; additional traded claims in the old span cannot help.

For $N=3$, normalization and one asset payoff $(0,1,2)$ priced at one imply $q_1=q_3=t$ and $q_2=1-2t$, with $0<t<1/2$. Every such measure prices the traded asset identically. A claim paying one only in state 2 has price $1-2t$, which spans $(0,1)$. Physical probabilities and a chosen entropy preference select a value, but asset calibration alone does not.

\section*{Remaining work}
An empirical selection rule needs preferences or ambiguity attitudes measured independently of the prices it is fitted to. Infinite-state closure, integrability and equivalence require additional conditions absent here. A utility-indifference price with wealth effects may be nonlinear in the claim and cannot be forced into a single probability measure. The benchmark establishes a valid selector and its domain of stability, while the general economic justification and identification agenda remains unresolved.

\section{Completion Estimate}
58\%

This is a post hoc, heuristic estimate for the full catalogue target.
Finite-state entropy minimization is proved to select a unique equivalent measure continuously under interior and rank conditions, and new-payoff distinguishability is characterized. Broader state spaces, independently identified economic preferences and nonlinear valuation domains remain outside the established benchmark.

\begin{thebibliography}{9}
\bibitem{review} Z. Chen and C. Kay (2026), Historical Developments in Probability Measures for Asset Pricing: From State Prices to Modern Pricing Kernels, arXiv:2605.27658. Primary PDF inspected for review identity and selection agenda. \url{https://arxiv.org/pdf/2605.27658}.
\end{thebibliography}

\end{document}
